\abstractch

随着大型语言模型（Large Language Model，LLM）能力的快速提升，以多智能体协作为核心的自动化工作流系统在复杂垂直领域任务中展现出广阔应用前景。然而，现有工作流自动生成方法面向电力系统等专业领域时仍存在两类关键问题：一方面，生成过程过度依赖 LLM 参数化知识，难以充分利用领域文档、设备关系和操作规程中的结构化知识；另一方面，子任务与智能体之间缺乏能力、工具和接口层面的量化匹配机制，容易造成角色分配不准确、工作流可解释性不足等问题。

针对上述问题，提出一种基于图检索增强生成（Graph Retrieval-Augmented Generation，GraphRAG）和三维量化匹配的智能体工作流自主生成方法。首先，设计面向电力领域的多模态知识图谱构建流程，通过实体关系抽取、视觉语言模型描述和 BGE-M3 向量嵌入，将电力领域文本、图像与流程图融合为统一图结构，并写入 Neo4j 图数据库。其次，设计基于子图扩展的 GraphRAG 检索方法，以语义向量检索获得种子节点，以广度优先遍历补充多跳邻域关系，形成结构化图谱上下文。进一步，构造子任务-智能体混合匹配算法，从能力嵌入相似度、工具覆盖率、输入输出兼容性三个维度建立加权评分函数，实现分解子任务到专用智能体的可解释映射。最后，设计工作流有向无环图（Directed Acyclic Graph，DAG）编排与三层迭代验证-修复机制，支持拓扑排序、循环检测、接口兼容性修复以及 LangGraph 格式导出。

围绕电网故障排查（OPS）、电力知识问答（KQA）和电力数据分析（DA）三类场景构建包含 30 条任务的分层基准测试集，并结合 PowerGridQA 问答数据和 OPSD 开放电力时序数据开展实验。实验结果表明：（1）在工作流可执行率上，所提方法和基线方法均达到 100\%，验证了 DAG 验证机制的有效性；（2）三维量化匹配算法是角色分配准确率（RAA）的决定性因素，相比无匹配的随机分配基线在全场景均值上提升 36.9 个百分点（0.056 $\rightarrow$ 0.425），且在三类场景中均稳定成立；（3）经 BGE-M3 语义化改造的步骤完整性指标将覆盖率从词汇版的 0.18 量级提升至 0.77 量级，更准确地反映了中文电力工作流的步骤覆盖；（4）BFS 子图扩展在当前 4{,}337 节点的电力知识图谱上未带来稳定增益，反而引入噪声节点扰动下游任务分解，提示后续应研究基于查询难度的自适应跳数控制方法。研究结果验证了基于图谱语义检索与三维量化匹配在专业领域工作流自动生成任务上的有效性。

\vspace{5pt}

\keywordsch 多模态知识图谱；图检索增强生成；智能体工作流；子任务-智能体匹配；电力系统
